% Part: sets-functions-relations
% Chapter: relations
% Section: orders

\documentclass[../../../include/open-logic-section]{subfiles}

\begin{document}

\olfileid{sfr}{rel}{ord}
\olsection{ক্রমসম্পর্ক}

\begin{explain}
আমাদের অনেক তুলনাতেই কিছু বস্তুকে কোনো বিশেষ দিক থেকে অন্য বস্তুর ``চেয়ে ছোট'', ``সমান'' বা ``চেয়ে বড়'' বলা হয়। এগুলি \emph{ক্রম}সম্পর্কের উদাহরণ। কিন্তু ক্রমসম্পর্কের নানা ধরন আছে। যেমন, কোনো কোনো ধরনের ক্ষেত্রে যেকোনো দুটি বস্তুকে তুলনা করা যায়, অন্যগুলির ক্ষেত্রে তা আবশ্যক নয়। কোনোটি অভিন্নতাকে অন্তর্ভুক্ত করে (যেমন~$\le$), আবার কোনোটি বাদ দেয় (যেমন~$<$)। তাই এখানে একটি শ্রেণিবিভাগ কাজে আসবে।
\end{explain}

\begin{defn}[প্রাক্‌ক্রম]
যে সম্পর্ক প্রতিবিম্ব ধর্মবিশিষ্ট এবং পরিযায়ী, তাকে \emph{প্রাক্‌ক্রম} বলা হয়।
\end{defn}

\begin{defn}[আংশিক ক্রম]
যে প্রাক্‌ক্রম বিপ্রতিসমও, তাকে \emph{আংশিক ক্রম} বলা হয়।
\end{defn}

\begin{defn}[রৈখিক ক্রম]\ollabel{def:linearorder}
যে আংশিক ক্রম সংযুক্তও, তাকে \emph{পূর্ণ ক্রম} বা \emph{রৈখিক ক্রম} বলা হয়।
\end{defn}

\begin{ex}
প্রত্যেক রৈখিক ক্রম আংশিক ক্রমও, এবং প্রত্যেক আংশিক ক্রম প্রাক্‌ক্রমও; কিন্তু বিপরীত দিকের দাবিগুলি সত্য নয়। $A$-এর উপর সার্বিক সম্পর্ক একটি প্রাক্‌ক্রম, কারণ তা প্রতিবিম্ব ধর্মবিশিষ্ট এবং পরিযায়ী। কিন্তু $A$-তে একটির বেশি !!{element} থাকলে সার্বিক সম্পর্ক বিপ্রতিসম নয়, তাই আংশিক ক্রমও নয়।
\end{ex}

\begin{ex}
$\Bin^*$-এর উপর \emph{চেয়ে দীর্ঘ নয়} সম্পর্ক $\preccurlyeq$ বিবেচনা করো: $x \preccurlyeq y$ যদি এবং কেবল যদি $\len{x} \le \len{y}$। এটি প্রাক্‌ক্রম (প্রতিবিম্ব ধর্মবিশিষ্ট এবং পরিযায়ী), এমনকি সংযুক্তও; কিন্তু আংশিক ক্রম নয়, কারণ এটি বিপ্রতিসম নয়। যেমন, $01 \preccurlyeq 10$ এবং $10 \preccurlyeq 01$, কিন্তু $01 \neq 10$।
\end{ex}

\begin{ex}
সেটের একটি সেটের উপর $\subseteq$ সম্পর্কটি একটি গুরুত্বপূর্ণ আংশিক ক্রম। সাধারণভাবে এটি রৈখিক ক্রম নয়; কারণ $a \neq b$ হলে $\Pow{\{a, b\}} = \{\emptyset, \{a\}, \{b\}, \{a,b\}\}$ বিবেচনা করে আমরা দেখি $\{a\} \nsubseteq \{b\}$ এবং $\{a\} \neq \{b\}$ এবং $\{b\} \nsubseteq \{a\}$।
\end{ex}

\begin{ex}
\emph{নিঃশেষে বিভাজ্যতার} সম্পর্ক থেকে আমরা এমন একটি আংশিক ক্রম পাই যা রৈখিক ক্রম নয়। পূর্ণসংখ্যা $n$ ও~$m$-এর জন্য $n \mid m$ লিখে আমরা বোঝাই যে $n$ দিয়ে $m$ নিঃশেষে বিভাজ্য, অর্থাৎ এমন একটি পূর্ণসংখ্যা~$k$ আছে যাতে $m = kn$ হয়। $\Nat$-এর উপর এটি আংশিক ক্রম, কিন্তু রৈখিক ক্রম নয়: যেমন, $2 \nmid 3$ এবং $3 \nmid 2$-ও। $\Int$-এর উপর সম্পর্ক হিসেবে বিবেচনা করলে বিভাজ্যতা কেবল প্রাক্‌ক্রম, কারণ তা বিপ্রতিসম নয়: $1 \mid -1$ এবং $-1 \mid 1$, কিন্তু $1 \neq -1$।
\end{ex}

\begin{ex}
অনুক্রমের সেট~$A^*$-এর উপর \emph{প্রসারণ} সম্পর্কটি এইরকম: $s \sqsubseteq s'$ যদি এবং কেবল যদি $s = \emptyseq$ (শূন্য অনুক্রম), $s = s'$, অথবা $s = \tuple{s_1, \dots, s_n}$ এবং $s' = \tuple{s_1, \dots, s_n, s_{n+1}, \dots, s_m}$ হয়। $s \sqsubseteq s'$ হলে আমরা $s$-কে $s'$-এর \emph{প্রারম্ভিক খণ্ড}ও বলি। $A^*$-এর উপর প্রসারণ সম্পর্ক আংশিক ক্রম, কিন্তু রৈখিক ক্রম নয়; যেমন, $a \neq b$ হলে $ab \not\sqsubseteq ba$ এবং $ba \not\sqsubseteq ab$।
\end{ex}

\begin{defn}[কঠোর ক্রম]
\emph{কঠোর ক্রম} হল এমন সম্পর্ক যা আত্মসম্পর্কহীন, একমুখী এবং পরিযায়ী।
\end{defn}

\begin{defn}[কঠোর রৈখিক ক্রম]\ollabel{def:strictlinearorder}
যে কঠোর ক্রম সংযুক্তও, তাকে \emph{কঠোর পূর্ণ ক্রম} বা \emph{কঠোর রৈখিক ক্রম} বলা হয়।
\end{defn}

\begin{ex}
$\le$ হল কঠোর রৈখিক ক্রম~$<$-এর অনুরূপ রৈখিক ক্রম। $\subseteq$ হল কঠোর ক্রম~$\subsetneq$-এর অনুরূপ আংশিক ক্রম।
\end{ex}

$A$-এর উপর যেকোনো কঠোর ক্রম $R$-কে কর্ণ $\Id{A}$ যোগ করে, অর্থাৎ সব~$\tuple{x, x}$ যুগল যোগ করে, আংশিক ক্রমে রূপান্তরিত করা যায়। (একে $R$-এর \emph{প্রতিবিম্ব আবরণ} বলা হয়।) বিপরীতভাবে, একটি আংশিক ক্রম থেকে শুরু করে $\Id{A}$ বাদ দিলে কঠোর ক্রম পাওয়া যায়। পরের দুটি ফল এই কথাগুলিই সুনির্দিষ্টভাবে প্রকাশ করে।

\begin{prop}\ollabel{prop:stricttopartial}
যদি $R$, $A$-এর উপর কঠোর ক্রম হয়, তবে $R^+ = R \cup \Id{A}$ আংশিক ক্রম। আরও, $R$ কঠোর রৈখিক ক্রম হলে $R^+$ রৈখিক ক্রম।
\end{prop}

\begin{proof}
ধরা যাক, $R$ কঠোর ক্রম; অর্থাৎ $R \subseteq A^2$ এবং $R$ আত্মসম্পর্কহীন, একমুখী ও পরিযায়ী। ধরি $R^+ = R \cup \Id{A}$। আমাদের দেখাতে হবে যে $R^+$ প্রতিবিম্ব ধর্মবিশিষ্ট, বিপ্রতিসম এবং পরিযায়ী।

$R^+$ যে প্রতিবিম্ব ধর্মবিশিষ্ট তা স্পষ্ট, কারণ সকল $x \in A$-এর জন্য $\tuple{x, x} \in \Id{A} \subseteq R^+$।

$R^+$ বিপ্রতিসম দেখাতে স্ববিরোধে পৌঁছানোর উদ্দেশ্যে ধরা যাক $R^+xy$ এবং $R^+yx$, কিন্তু $x \neq y$। যেহেতু $\tuple{x,y} \in R \cup \Id{A}$, অথচ $\tuple{x, y} \notin \Id{A}$, তাই অবশ্যই $\tuple{x, y} \in R$, অর্থাৎ $Rxy$। একইভাবে,~$Ryx$। কিন্তু এটি $R$ একমুখী হওয়ার অনুমানের বিরোধী।

পরিযায়িতা প্রতিষ্ঠা করতে ধরা যাক $R^+xy$ এবং $R^+yz$। যদি $\tuple{x, y} \in R$ এবং $\tuple{y,z} \in R$ উভয়ই হয়, তবে $R$~পরিযায়ী বলে $\tuple{x, z} \in R$। অন্যথায়, হয় $\tuple{x, y} \in \Id{A}$, অর্থাৎ $x = y$, নয়তো $\tuple{y, z} \in \Id{A}$, অর্থাৎ $y = z$। প্রথম ক্ষেত্রে অনুমান অনুসারে $R^+yz$, এবং $x = y$; তাই $R^+xz$। দ্বিতীয় ক্ষেত্রেও একইরকম। উভয় ক্ষেত্রেই $R^+xz$; সুতরাং $R^+$-ও পরিযায়ী।

``আরও'' অংশটির জন্য ধরা যাক $R$ সংযুক্তও। তাহলে সকল $x \neq y$-এর জন্য হয় $Rxy$, নয়তো~$Ryx$; অর্থাৎ হয় $\tuple{x, y} \in R$, নয়তো $\tuple{y, x} \in R$। যেহেতু $R \subseteq R^+$, তাই $R^+$-এর ক্ষেত্রেও এটি সত্য থাকে; সুতরাং $R^+$-ও সংযুক্ত।
\end{proof}

\begin{prop}\ollabel{prop:partialtostrict}
যদি $R$, $A$-এর উপর আংশিক ক্রম হয়, তবে $R^- = R \setminus \Id{A}$ কঠোর ক্রম। আরও, $R$ রৈখিক ক্রম হলে $R^-$ কঠোর রৈখিক ক্রম।
\end{prop}

\begin{proof}
এটি অনুশীলন হিসেবে রাখা হল।
\end{proof}

\begin{prob}
\olref[sfr][rel][ord]{prop:partialtostrict}-এর একটি প্রমাণ দাও।
\end{prob}

পরের সহজ ফলটি প্রতিষ্ঠা করে যে কঠোর রৈখিক ক্রম সদস্যভিত্তিক সমতার নীতির অনুরূপ একটি ধর্ম মেনে চলে:

\begin{prop}\ollabel{prop:extensionality-strictlinearorders}
যদি $<$, $A$-এর উপর কঠোর রৈখিক ক্রম হয়, তবে:
\[
  (\forall a, b \in A)((\forall x \in A)(x < a \liff x < b) \lif a = b).
\]
\end{prop}

\begin{proof}
ধরা যাক $(\forall x \in A)(x < a \liff x < b)$। যদি $a < b$ হয়, তবে $a < a$ হয়, যা $<$ আত্মসম্পর্কহীন হওয়ার বিরোধী; তাই $a \nless b$। ঠিক একইভাবে $b \nless a$। সুতরাং $a = b$, কারণ $<$ সংযুক্ত।
\end{proof}

\end{document}
